Comparing TCP congestion control and active queue management in ns-3 is highly sensitive to topology, RTT, ECN configuration, queue discipline defaults, flow mix, and random seeds, yet many simulation studies report only a handful of manually executed scenarios. This paper presents a reproducible ns-3 benchmark workflow for TCP/AQM studies and demonstrates it through baseline single-flow and mixed-flow comparisons of Linux Reno, CUBIC, and DCTCP with CoDel, FQ-CoDel, PIE, and RED. The workflow ships as a small contrib module that loads YAML or JSON experiment and topology descriptions, expands matrices over seeded random-run values, and preserves command lines, environment metadata, and raw traces per run alongside summary tables and vector figures. The benchmark scenario is derived from the in-tree TCP validation example, inheriting a known-good dumbbell topology and validation cases without modifying the upstream example. The case study reports throughput, fairness, RTT, queueing delay, drops, and ECN marks across base RTT, ECN, and competing-flow settings, distinguishing conclusions that hold across configurations from those that depend on specific assumptions. The artifact is designed to be rerun, audited, and extended by ns-3 users, and is available on the \texttt{icns3-summer-2026-benchmark} branch of our ns-3 fork.\footnote{\url{https://github.com/Haoyu2/ns3/tree/icns3-summer-2026-benchmark}}
